%% kapa: exam=1 round=21 year=2010 period=1 form=A subject=economics subject_ko=경제원론 seq=18 number=58
%% source: Q-Net 감정평가사 기출문제 (cst00309) artlSeq=5024181 file=extracted/1008/5024181/2010년 제21회 감정평가사 1차시험 기출문제/1교시.hwp
%% conversion: kordoc 4.24.51 -> tools/gen1.py (auto); review=proofread-sonnet; flags=none
\begin{question}{58}
甲은 열등재인 X재와 정상재인 Y재만을 소비한다. 소득과 Y재의 가격이 일정할 때, X재의 가격이 하락하자 X재의 소비량이 감소하였다. 이 경우 옳은 설명을 모두 고른 것은?
\begin{framedtext}
ㄱ. X재는 기펜재이다. \par
ㄴ. 가격소비곡선은 우상향한다. \par
ㄷ. X재에 대한 대체효과의 절대값이 소득효과의 절대값보다 작다. \par
ㄹ. X재에 대한 대체효과와 소득효과는 같은 방향이다. \par
\end{framedtext}
\begin{choices}
  \item ㄱ, ㄴ
  \item ㄱ, ㄷ
  \item ㄴ, ㄹ
  \item ㄱ, ㄴ, ㄷ
  \item ㄴ, ㄷ, ㄹ
\end{choices}
\end{question}
\answer{2}
